\documentclass[11pt,a4paper]{article}
\usepackage[utf8]{inputenc}
\usepackage[T1]{fontenc}
\usepackage[portuguese]{babel}
\usepackage{xcolor}
\usepackage{hyperref}
\usepackage{geometry}
\usepackage{listings}
\usepackage{tcolorbox}
\geometry{margin=2.5cm}
\definecolor{primary}{RGB}{41,128,185}
\definecolor{secondary}{RGB}{44,62,80}
\definecolor{codebg}{RGB}{240,240,240}
\lstset{backgroundcolor=\color{codebg},basicstyle=\ttfamily\small,breaklines=true,frame=single}
\title{\color{secondary}\Huge\textbf{Recursos de GKE, Ingress, TLS e Canary Deployment}}
\author{\color{primary}DevOps com Kubernetes -- 2026}
\date{}
\begin{document}
\maketitle


\section*{\color{primary}O que você aprenderá nesta página}
\begin{itemize}
\item Publicar aplicações com Ingress ou Gateway
\item Configurar HTTPS com TLS e cert-manager
\item Usar labels, annotations e Canary Deployment para controlar tráfego
\end{itemize}

Depois que a aplicação está no cluster e o pipeline consegue publicar novas versões, o próximo passo é lidar com tráfego externo, segurança de transporte e estratégias de lançamento. Em cloud, esses temas aparecem juntos: o provedor oferece balanceadores, o Kubernetes oferece objetos declarativos e ferramentas como cert-manager automatizam certificados.

\section*{\color{primary}Ingress e Ingress Controller}

\texttt{Ingress} é um recurso que descreve regras HTTP e HTTPS para chegar aos Services internos. Ele não funciona sozinho: é necessário um \texttt{Ingress Controller}, como NGINX Ingress Controller, Traefik ou controlador do provedor de nuvem.

\begin{lstlisting}[language=bash]
apiVersion: networking.k8s.io/v1
kind: Ingress
metadata:
  name: app-ingress
spec:
  rules:
    - host: app.example.com
      http:
        paths:
          - path: /
            pathType: Prefix
            backend:
              service:
                name: app-svc
                port:
                  number: 80
\end{lstlisting}

\section*{\color{primary}TLS e cert-manager}

Para HTTPS, o Ingress precisa de um certificado TLS. Em estudos, esse certificado pode ser criado manualmente. Em produção, é comum usar \texttt{cert-manager} para emitir e renovar certificados automaticamente.

\begin{lstlisting}[language=bash]
apiVersion: cert-manager.io/v1
kind: Certificate
metadata:
  name: app-tls
spec:
  secretName: app-tls-secret
  dnsNames:
    - app.example.com
  issuerRef:
    name: letsencrypt-prod
    kind: ClusterIssuer
\end{lstlisting}

\section*{\color{primary}Labels e annotations}

Labels são usadas para seleção e organização. Annotations armazenam metadados usados por controladores, como Ingress Controller, cert-manager e ferramentas de observabilidade. Use labels para aquilo que será selecionado; use annotations para instruções auxiliares interpretadas por controladores externos.

\section*{\color{primary}Canary Deployment}

Canary Deployment libera uma nova versão para uma parte pequena do tráfego antes de substituir tudo. Com Ingress Controller, isso pode ser feito por annotations específicas; com ferramentas como Argo Rollouts, pode ser feito de forma mais estruturada e observável. A ideia é reduzir risco antes de ampliar o tráfego para todos os usuários.

\section*{\color{primary}Fechamento do capítulo}

Nesta parte, a aplicação deixou de ser apenas um Pod em execução. Ela passou a ter configuração externa, persistência, rede estável, publicação HTTP/HTTPS, pipeline e observabilidade. Esse é o caminho para transformar uma aplicação básica em uma aplicação preparada para cloud.

\end{document}
